%%%%%%%%%%%%%%%%%%%%%%%%%%%%%%%%%%%%%%%%%
% XXV Simposio Chileno de Física, LaTeX Template, Versión 1 (8 may 2026)
% vmunoz@fisica.ciencias.uchile.cl
% Basado en versión previa de
% fabian.torres@ufrontera.cl    (7 Junio 2024)
% daniel.salinas@sansano.usm.cl (12 Septiembre 2022)
%%%%%%%%%%%%%%%%%%%%%%%%%%%%%%%%%%%%%%%%%
% Método de compilación:
%
%				LaTeX->PDF
%
%%%%%%%%%%%%%%%%%%%%%%%%%%%%%%%%%%%%%%%%%
%
%                                                    INSTRUCCIONES SOBRE EL DOCUMENTO
%
% - Subrayar al expositor del trabajo,
% -	Borrar el texto en rojo y este comentario
% -	No se debe modificar los tamaños de letra.
% -	El título del trabajo no debe superar 2 líneas, o 30 palabras o 140 caracteres incluyendo espacios.
% -	Preferentemente, incluya sólo el correo electrónico del expositor del trabajo.
% -	Incluya un máximo de 5 autores. En caso de tener más de 5 autores, seleccione los principales incluyendo al expositor, y agregue et al %   al listado para incluir al resto.
%
%-----------------------------------------

%----------------------------------------------------------------------------------------
%	CONFIGURACIÓN DEL DOCUMENTO
%----------------------------------------------------------------------------------------

\documentclass[11pt,spanish]{article}

\usepackage[spanish]{babel}
\usepackage{epsfig,graphicx}
\usepackage{amsmath, amsthm, amsfonts}
\usepackage[hmarginratio=6:5,top=32mm,columnsep=11mm,textwidth=178mm]{geometry}
\usepackage[makeroom]{cancel} %FT - para incluir simplificaciones en ecuaciones
\usepackage{multicol} %FT - para generar varias columnas en una seccion del documento
\usepackage[hang,small,labelfont=bf,up,textfont=it,up]{caption} % Custom captions under/above floats in tables or figures
\usepackage{booktabs} % Horizontal rules in tables
\usepackage{float} % Required for tables and figures in the multi-column environment - they need to be placed in specific locations with the [H] (e.g. \begin{table}[H])
\usepackage{hyperref} % For hyperlinks in the PDF
\usepackage{paralist} % Used for the compact item environment which makes bullet points with less space between them

\addcontentsline{toc}{section}{Title of the section}
\usepackage{titlesec} % Allows customization of titles
\usepackage{fancyhdr} % Headers and footers


\pagestyle{fancy} % All pages have headers and footers
\fancyhead{} % Blank out the default header
\fancyfoot{} % Blank out the default footer
\renewcommand{\headrulewidth}{3pt}
\renewcommand{\footrulewidth}{0pt}
%-----------------------------------------------------------------------------------------
% AGREGUE AQUÍ UNA DE LAS ÁREAS DEFINIDAS EN LA PÁGINA WEB DEL SIMPOSIO www.sochifi2014.udec.cl
%-----------------------------------------------------------------------------------------


\fancyhead[L]{\small 	Física mas allá \\ del Modelo Estándar} % Custom header text
\fancyhead[R]{{\bf XXV Simposio Chileno de F\'isica} \\  23 al 26 de
  noviembre de 2026 \\  Santiago, Chile} % Custom header text
\fancyfoot[RO,L]{\thepage} % Custom footer text

\linespread{1} %

%----------------------------------------------------------------------------------------
%	TÍTULO Y AUTORES. SUBRAYAR AL EXPOSITOR
%----------------------------------------------------------------------------------------

\title{\vspace{-5mm}\fontsize{16pt}{10pt}\selectfont \textbf{Efecto Sommerfield en modelos de materia oscura con multipletes de SU(2)_L}} % Article title
\vspace{-3mm}
\author{
\normalsize Marcela González$^{1*}$, \underline{Jovanny Jara}$^{1 \dagger}$
 \\ \vspace{-1mm}
\small  $^1$Instituto de Física y Astronomía, Universidad de Valparaíso, Valparaíso \\ %Institución 1
\small $*$marcela.gonzalezpi@uv.cl, $\dagger$jovanny.jara@estudiantes.uv.cl}
\date{\today}
\pagestyle{empty}
%----------------------------------------------------------------------------------------

\newcommand{\seccion}[1]{\noindent {\Large\textbf{#1}}\\}

\begin{document}
\maketitle
\thispagestyle{fancy} % All pages have headers and footers
\pagenumbering{gobble}
%----------------------------------------------------------------------------------------
%	CONTENIDO
%----------------------------------------------------------------------------------------

\seccion{Resumen}
\noindent
La naturaleza de la materia oscura (DM) es uno de los problemas centrales de la cosmología y la física de partículas. Diversas observaciones astronómicas y cosmológicas sugieren que la mayor parte de la materia en el universo corresponde a DM, pero su composición y propiedades siguen siendo desconocidas.\\
Una posibilidad es que la DM corresponda a una partícula masiva estable introducida en modelos de nueva física más allá del Modelo Estándar (SM), cuya densidad reliquia se determina con la sección eficaz de aniquilación. Al considerar modelos donde las partículas de DM corresponden a multipletes no triviales del grupo de gauge electrodébil $SU(2)_L$ entre las cuales existe una interacción a largo alcance, la sección eficaz de aniquilación puede verse significativamente aumentada en el límite no relativista, fenómeno que se conoce como Sommerfield Enhancement, relevante para la fenomenología de DM.\\
% En esta sección obligatoria, escriba un resumen e introducción del trabajo que presenta al Simposio.
% El texto de la contribución puede estar en castellano o en inglés.
% No se extienda más allá de una página tamaño carta y no modifique
% el tipo de letra. Subraye el nombre del autor que presenta el trabajo
% (sea oral o poster), tal como se indica más arriba, e incluya su
% correo electrónico después de las afiliaciones. Incluya un máximo
% de 5 autores. En caso de tener más de 5 autores, seleccione los
% principales incluyendo al expositor, y agregue et al para incluir al
% resto. Use como índices los símbolos $*$, $\dag$, $\ddag$, o $**$.
% Puede dividir o no el texto en secciones, aparte de esta primera
% sección que es obligatoria. Introduzca las referencias manualmente y
% use números entre corchetes $[1]$. Dos referencias juntas escríbalas
% $[1, 2]$. Agregue la lista de referencias al final, usando el mismo
% formato de las secciones. Envíe su contribución mediante el
% formulario dispuesto en la página web del encuentro. Para consultas
% escriba al Comité Organizador, al correo: {\tt organizacion@sochifi2026.cl}.\\


%Incluya ecuaciones, tablas y figuras si es necesario.


\seccion{Desarrollo}
\noindent
El factor de Sommerfield $S$ se define como
\begin{equation}
	S = \frac{|\psi(0)|^2}{|\psi_0(0)|^2},
\end{equation}
donde $\psi(0)$ es la función de onda no relativista considerando el potencial de interacción y $\psi_0(0)$ es la función de onda libre [1]. Así, el problema se trabaja resolviendo la ecuación de Schrödinger para $\psi(0)$ y $\psi_0(0)$. \\
Al introducir modelos de multipletes de $SU(2)_L$ se trabajan con potenciales con forma matricial y el problema se describe mediante un sistema de ecuaciones de Schrödinger acopladas de la forma [2]
\begin{equation}
	\left(-\frac{1}{m_\xi}\nabla^2 + V_{ij}(r)\right)\psi_j = \psi_i.
\end{equation}
% Las tablas pueden construirse en el archivo directamente o insertarse como cuadro de texto o imagen. Las figuras, por su parte, deben insertarse como imagen. También puede usar dos columnas si lo necesita.

% \begin{multicols}{2}
% 	\begin{table}[H]
% 		\centering
% 		\begin{tabular}{|l|l|l|}
% 			\hline
% 			Nombre  & Apellido & Nota   \\ \hline
% 			Diego   & Kosta    & $6.0$  \\ \hline
% 			Antonio & Grassman & $5.5 $ \\ \hline
% 		\end{tabular}\\
% 		\vspace{0.3cm}
% 		Tabla I. Ejemplo
% 	\end{table}

% 	\columnbreak
% 	\begin{figure}[H]
% 		\centerline{\includegraphics[width=3cm]{sochifilogo.png}}
% 	\end{figure}
% \end{multicols}


%------------------------------------------------

% \noindent \textbf{Agradecimientos:} Incluya los agradecimientos antes de las referencias.\\


%------------------------------------------------

% Escriba aqu\'i s\'olo las referencias m\'as importantes del trabajo



\noindent {\large\textbf{Referencias}}


\noindent
[1] Jiménez Arranz, Ó. The role of the Sommerfeld enhancement in dark matter physics. (2018).\newline
[2] M. Cirelli, A. Strumia and M. Tamburini, Cosmology and Astrophysics of Minimal Dark
Matter, Nucl. Phys. B 787 (2007) 152 [hep-ph/0706.4071].











%----------------------------------------------------------------------------------------



\end{document}
